\documentclass{article}

% Packages
\usepackage{fancyhdr}
\usepackage{extramarks}
\usepackage{amsmath}
\usepackage{amsthm}
\usepackage{amsfonts}
\usepackage{tikz}
\usepackage[plain]{algorithm}
\usepackage{algpseudocode}
\usepackage{enumerate}

\usetikzlibrary{automata,positioning}

% Document Layout
\topmargin=-0.45in
\evensidemargin=0in
\oddsidemargin=0in
\textwidth=6.5in
\textheight=9.0in
\headsep=0.25in
\linespread{1.1}

% Page Style
\pagestyle{fancy}
\lhead{\hmwkAuthorName}
\chead{\hmwkClass:\ \hmwkTitle}
\rhead{Section \hmwkSection, \firstxmark}
\lfoot{\lastxmark}
\cfoot{\thepage}
\renewcommand\headrulewidth{0.4pt}
\renewcommand\footrulewidth{0.4pt}

% Paragraph Settings
\setlength\parindent{0pt}
\setlength{\parskip}{5pt}

% Section Management
\newcommand{\hmwkSection}{A} % Current section (A, B, or C) - update manually

% Problem Header Management
\newcommand{\enterProblemHeader}[1]{
  \nobreak\extramarks{}{Problem \arabic{#1} continued on next page\ldots}\nobreak{}
  \nobreak\extramarks{Problem \arabic{#1} (continued)}{Problem \arabic{#1} continued on next page\ldots}\nobreak{}
}

\newcommand{\exitProblemHeader}[1]{
  \nobreak\extramarks{Problem \arabic{#1} (continued)}{Problem \arabic{#1} continued on next page\ldots}\nobreak{}
  \stepcounter{#1}
  \nobreak\extramarks{Problem \arabic{#1}}{}\nobreak{}
}

% Counters
\setcounter{secnumdepth}{0}
\newcounter{partCounter}
\newcounter{homeworkProblemCounter}
\setcounter{homeworkProblemCounter}{1}
\nobreak\extramarks{Problem \arabic{homeworkProblemCounter}}{}\nobreak{}

% Homework Problem Environment
% Optional argument adjusts problem counter for non-sequential problems
\newenvironment{homeworkProblem}[1][-1]{
  \ifnum#1>0
    \setcounter{homeworkProblemCounter}{#1}
  \fi
  \section{Problem \arabic{homeworkProblemCounter}}
  \setcounter{partCounter}{1}
  \enterProblemHeader{homeworkProblemCounter}
}{
  \exitProblemHeader{homeworkProblemCounter}
}

% Assignment Details
\newcommand{\hmwkTitle}{Sheet\ \#2}
\newcommand{\hmwkDueDate}{February 12, 2026}
\newcommand{\hmwkClass}{C8.4 Probabilistic Combinatorics}
\newcommand{\hmwkClassInstructor}{Professor P. Balister}
\newcommand{\hmwkAuthorName}{\textbf{Ray Tsai}}

% Title Page
\title{
  \vspace{2in}
  \textmd{\textbf{\hmwkClass:\ \hmwkTitle}}\\
  \normalsize\vspace{0.1in}\small{Due\ on\ \hmwkDueDate\ at 12:00pm}\\
  \vspace{0.1in}\large{\textit{\hmwkClassInstructor}} \\
  \vspace{3in}
}

\author{\hmwkAuthorName}
\date{}

% Part Command
\renewcommand{\part}[1]{\textbf{\large Part \Alph{partCounter}}\stepcounter{partCounter}\\}

% Mathematical Commands
% Algorithms
\newcommand{\alg}[1]{\textsc{\bfseries \footnotesize #1}}

% Calculus
\newcommand{\deriv}[1]{\frac{\mathrm{d}}{\mathrm{d}x} (#1)}
\newcommand{\pderiv}[2]{\frac{\partial}{\partial #1} (#2)}
\newcommand{\dx}{\mathrm{d}x}

% Probability and Statistics
\newcommand{\Var}{\mathrm{Var}}
\newcommand{\Cov}{\mathrm{Cov}}
\newcommand{\Bias}{\mathrm{Bias}}
\newcommand*{\prob}{\mathbb{P}}
\newcommand*{\E}{\mathbb{E}}

% Number Sets
\newcommand*{\Z}{\mathbb{Z}}
\newcommand*{\Q}{\mathbb{Q}}
\newcommand*{\R}{\mathbb{R}}
\newcommand*{\C}{\mathbb{C}}
\newcommand*{\N}{\mathbb{N}}

\begin{document}

\maketitle
\pagebreak

\begin{homeworkProblem}
  Show that all (finite) graphs of the following types are strictly balanced: complete graphs, cycles, trees, complete bipartite graphs, connected regular graphs.

  Give (with proof) an example of a graph that is balanced but not strictly balanced.

  \begin{proof}
    We first note that for any subgraph $H$ of $G$ with vertex set $U \subseteq V(G)$, we have $d(H) \leq d(G[U])$. Thus it suffices to show that for $d(G[U]) < d(G)$ for all $U \subseteq V(G)$.

    \textbf{Complete graphs:} Let $G$ be a complete graph with $n$ vertices. Note that all induced subgraphs of $G$ are also complete graphs. But then the function $f(n) = \binom{n}{2}/n$ is strictly increasing on $n$, so $d(G[U]) < d(G)$ for all $U \subseteq V(G)$.

    \textbf{Cycles:} Let $G$ be a cycle with $n$ vertices. Note that for all strict subsets $U \subset V(G)$, the induced subgraph $G[U]$ is a collection of disjoint paths, which has density
    \[
      d(G[U]) \leq \frac{|U| - 1}{|U|} < 1 = d(G).
    \]

    \textbf{Trees:} Let $G$ be a tree with $n$ vertices. Note that all induced subgraphs of $G$ are forests. Thus for $U \subseteq V(G)$, the density of of the induced subgraph is
    \[
      d(G[U]) \leq \frac{|U| - 1}{|U|} = 1 - \frac{1}{|U|},
    \]
    which is strictly increasing on $|U| \geq 1$. Hence, $d(G[U]) < d(G)$ for all $U \subseteq V(G)$.

    \textbf{Complete bipartite graphs:} Let $G$ be a complete bipartite graph with $n$ vertices with parts $A$ and $B$. Note that all induced subgraphs of $G$ are also complete bipartite graphs. Let $U \subseteq V(G)$ and let $a = |A \cap U|$ and $b = |B \cap U|$. By AM-GM
    \[
      d(G[U]) = \frac{ab}{a + b} \leq \frac{(a + b)^2}{4(a + b)} = \frac{|U|}{4},
    \]
    which is strictly increasing on $|U|$. Hence, $d(G[U]) < d(G)$ for all $U \subseteq V(G)$.

    \textbf{Connected regular graphs:} Let $G$ be a connected regular graph with $n$ vertices with each vertex having degree $k$. Let $H$ be a strict subgraph of $G$. Note that $H$ cannot be $k$-regular, otherwise $H$ would be it's own component. Thus any $H$ must contain some vertex with degree less than $k$. But then
    \[
      d(H) < \frac{k|V(H)|}{2|V(H)|} = \frac{k}{2} = d(G).
    \]

    \textbf{Balanced graphs:} Consider $G$ being a union of two cycles $C_1$ and $C_2$. We already showed that $C_1$ and $C_2$ are balanced. It now follows $d(G) = 1 = d(C_1)$ that $G$ is not strictly balanced.
  \end{proof}
\end{homeworkProblem}

\newpage

% To change sections, use: \renewcommand{\hmwkSection}{B}
% Example: Uncomment the line below when you start Section B
\renewcommand{\hmwkSection}{B}

\begin{homeworkProblem}
  Let $T$ be a tree with $k$ vertices. What is the threshold for $G(n, p)$ to contain a copy of $T$?
    
  Show that, for each fixed $k$, there is a function $p(n)$ such that the probability that $G(n, p(n))$ has a component of size exactly $k$ tends to 1 as $n \to \infty$.

  \begin{proof}
    By problem 1, we know that trees on $k$ vertices are balanced. Thus, it has threshold $n^{-1 - \frac{1}{k - 1}}$, by Theorem 2.5. Set $p = c/n$ for some constant $c < 1$. Let $X$ be the number of components of size $k$, $Y$ be the number of trees of size $k$, and $C$ be the number of cycles in $G(n, p)$. Note that
    \[
      \E[C] = \sum_{k = 3}^{n} \binom{n}{k} \cdot \frac{(k - 1)!}{2} \cdot p^k \leq \frac{1}{2}\sum_{k = 0}^{n} \frac{c^k}{k} \leq -\log(1 - c) = O(1).
    \]
    Thus $X = Y + O(1)$, almost surely. By Cayley's formula, there are $\binom{n}{k}k^{k - 2}$ trees on $k$ vertices in $K_n$. Thus,
    \[
      \E[X] = \binom{n}{k}k^{k - 2}p^{k - 1}(1 - p)^{\binom{k}{2} - k + 1}(1 - p)^{(n - k)k}.
    \]
    Since $\binom{n}{k} = \Theta(n^k)$ and $1 - p \approx e^{-p}$, we have
    \[
      \E[X] = \Theta(n^k) \cdot \Theta(n^{k - 1}) \cdot e^{-p(1 + o(1)kn)} = \Theta(n).
    \]
    For $S \subseteq V(G)$ with $|S| = k$, let $A_S$ be the event that $S$ induces an isolated tree. We note that
    \[
      \prob(A_S) = (1 - p)^{(n - k)k} \cdot \prob(S \text{ induces a tree})
    \] 
    Summing over all ordered pairs of sets of vertices $S, S'$ of size $k$,
    \[
      \E[X(X - 1)] = \sum_{S \neq S'} \prob(A_S \cap A_{S'})
    \]
    If $S$ and $S'$ intersect, then it is impossible for both $S$ and $S'$ to induce an isolated tree. Now suppose $S \cap S' = \emptyset$. Then
    \[
      \prob(A_S \cap A_{S'}) = (1 - p)^{2(n - k)k - k^2}\prob(S \text{ induces a tree}) \cdot \prob(S' \text{ induces a tree}). 
    \]
    But then
    \[
      \frac{\prob(A_S \cap A_{S'})}{\prob(A_S) \prob(A_{S'})} = \frac{(1 - p)^{2(n - k)k - k^2}}{((1 - p)^{(n - k)k})^2} = (1 - p)^{-k^2} \approx e^{-ck^2/n} \to 1,
    \]
    as $n \to \infty$. Thus,
    \[
      \E[X(X - 1)] \sim \sum_{S \neq S'} \prob(A_S)^2 = \sum_{S}\sum_{S'} \prob(A_S)\prob(A_{S'}) - \E[X] \sim \E[X]^2.
    \]
    It now follows from Chebyshev's inequality that
    \[
      \prob(Y = 0) \leq \prob(X = 0) \leq \prob(|X - \E[X]| \geq \E[X]) \leq \frac{\E[X^2]}{\E[X]^2} - 1 \to 0,
    \]
    as $n \to \infty$.
  \end{proof}
\end{homeworkProblem}

\newpage

\begin{homeworkProblem}
  What is the threshold for $G(n, p)$ to contain a cycle of length $k$, for fixed $k$?
    
  Find a threshold function for the property of containing a cycle. [\textit{Careful! It doesn't quite follow immediately from the first part.}]

  \begin{proof}
    By problem 1, cycles are balanced, so the threshold for cycles of length $k$ is $n^{-1}$. 

    Let $X$ be the expected number of cycles in $G(n, p)$. Then
    \[
      \E[X] = \sum_{k = 3}^n \binom{n}{k} \cdot \frac{(k - 1)!}{2} \cdot p^k = \frac{1}{2}\sum_{k = 3}^n \frac{(np)^k}{k} \sim \frac{1}{2}\left(1 + np + \frac{(np)^2}{2} - \log(1 - np)\right),
    \]
    as $n \to \infty$. Thus $\E[X] \to 0$ if $np \to 0$, and $\E[X] \to \infty$ if $np \to \infty$. This shows that $\prob(X > 0) \to 0$ for $np \to 0$. If $np \to \infty$, then $p$ passes the threshold for cycles of length $k$, so $\prob(X > 0) \to 1$.
  \end{proof}
\end{homeworkProblem}
\newpage

\begin{homeworkProblem}
  Show that it is possible to colour the edges of $K_n$ with $k = \lfloor 3\sqrt{n} \rfloor$ colours so that no triangle has all its edges the same colour.

  \begin{proof}
    Randomly color the edges of $K_n$ with $k$ colors. For triangle $i$ in $K_n$, let $A_i$ be the event that $i$ is monochromatic. Note that a triangle is adjacent to $d = 3(n - 3)$ other triangles in $K_n$, so $A_i$ is independent of the family of all but $d$ other $A_j$'s. Since 
    \[
      \prob(A_i) \leq \frac{1}{9n} \leq \frac{1}{e(3n - 8)},
    \]
    we have $\prob(\bigcap_{i = 1}^n A_i^c) > 0$ by LLL.
  \end{proof}
\end{homeworkProblem}

\newpage

\begin{homeworkProblem}
  Let $H = (V, E)$ be a hypergraph. Suppose the vertices are $k$-coloured uniformly at random; each $v \in V$ receives each colour with probability $1/k$, and the colours of different vertices are independent. For $e \in E$, let $A_e$ be the event that the hyperedge $e$ is monochromatic.
  
  Show that if $|e \cap f| \le 1$, then $A_e$ and $A_f$ are independent.
  
  Is it true that $A_e$ is independent of the collection $\{A_f : |e \cap f| \le 1\}$ of events?

  \begin{proof}
    If $e \cap f  = \emptyset$ then obviously $A_e$ and $A_f$ are independent. Now suppose $e \cap f = \{v\}$. Then the only vertex in $e$ that is affected by $A_f$ is $v$. But then the color of $v$ is independent of the event $A_f$. Thus, $A_e$ and $A_f$ are independent.

    However, $A_e$ is not necessarily independent of the collection $\{A_f : |e \cap f| \le 1\}$. Consider the cycle of length 4 $C_4$. If three of the edges are monochromatic, then the fourth edge must be monochromatic. 
  \end{proof}
\end{homeworkProblem}

\newpage

\begin{homeworkProblem}
  A $k$-SAT (more properly $k$-CNF) formula is an expression such as (for $k=3$)
  \[ (x_1 \text{ OR } x_4 \text{ OR } \overline{x_6}) \text{ AND } (\overline{x_1} \text{ OR } x_2 \text{ OR } x_5) \]
  \[ \text{AND } (\overline{x_2} \text{ OR } \overline{x_3} \text{ OR } \overline{x_4}) \text{ AND } (\overline{x_4} \text{ OR } \overline{x_5} \text{ OR } x_6), \]
  where the variables $x_i$ take values TRUE or FALSE, $\overline{x_i}$ means NOT $x_i$, and $k$ distinct variables or their negations are ORd together in each clause. A formula is satisfiable if there is an assignment of values to the variables making the expression true.
  
  \begin{enumerate}[(a)]
    \item Use the first-moment method to show that each $k$-SAT formula with fewer than $2^k$ clauses is satisfiable.
    \begin{proof}
      Suppose there are $m < 2^k$ clauses. Randomly assign values to each variable with probability $1/2$. Let $X$ be the number of clauses that are not satisfied. Then
      \[
        \prob(X \geq 1) = \E[X] = \sum_{i = 1}^m \prob(\text{clause $i$ is not satisfied}) =  \frac{m}{2^k} < 1.
      \]
      Thus, there exists an assignment of values to the variables that satisfies all the clauses.
    \end{proof}
    \item Use the Symmetric Local Lemma to show that each $k$-SAT formula in which no variable appears in more than $2^{k-2}/k$ clauses is satisfiable.
    \begin{proof}
      Randomly assign values to each variable with probability $1/2$. Let $A_i$ be the event that the $i$th clause is not satisfied. Since each clause contains $k$ variables and each variable appears in at most $2^{k-2}/k$ clauses, $A_i$ is dependent on at most $2^{k-2} - 1$ other $A_j$'s. But then $\prob(A_i) = 2^{-k} \leq \frac{1}{e(2^{k-2})}$. The result now follows from the Symmetric Local Lemma.
    \end{proof}
  \end{enumerate}
\end{homeworkProblem}

\newpage

\begin{homeworkProblem}
  Let $G = (V, E)$ be a graph with maximum degree $\Delta$, and let $V_1, V_2, \dots, V_s$ be a collection of disjoint subsets of $V$, each of size $k \ge 2e\Delta$. Show that $G$ has an independent set which contains a vertex from each set $V_i$. 

  \begin{proof}
    Pick a vertex $v_i$ from each $V_i$ uniformly at random, and let $I = \{v_1, v_2, \dots, v_s\}$. For $\{u, v\} \in E$, let $A_{uv}$ be the event that $\{u, v\} \subseteq I$. Note that the $A_{uv}$ is dependent to only $A_{u'v'}$ where $u, u'$ or $v, v'$ share the same set $V_i$. Thus, the $A_uv$ is independent of the family of all but at most $d = 2\Delta(2e\Delta) - 1$ other $A_{e'}$. But then $\prob(A_e) \leq 1/(2e\Delta)^2 = 1/e(d + 1)$. The result now follows from the Symmetric Local Lemma.
  \end{proof}
\end{homeworkProblem}

\newpage

\begin{homeworkProblem}
  Let $G = (V, E)$ be a graph, and suppose that for each $v \in V$ there is a list $S(v)$ of at least $2er$ colours, where $r$ is a positive integer. Suppose also that for each $v \in V$ and each $c \in S(v)$, there are at most $r$ neighbours $u$ of $v$ such that $c \in S(u)$.
  
  Prove that there is a proper colouring of $G$ under which each vertex $v$ receives a colour from its list $S(v)$.

  \begin{proof}
    For each vertex $v \in V$, randomly assign it a color from $S(v)$ uniformly at random. For $e \in E$, let $A_e$ be the event that $e$ is monochromatic. Note that each edge $e$ is incident to at most $2r - 1$ other edges such that $e \cup e'$ could be monochromatic. Thus $A_e$ is independent of the family of all but at most $d = 2r - 1$ other $A_{e'}$. Since $\prob(A_e) \leq 1/2er = 1/ed$, the result now follows from the Symmetric Local Lemma.
  \end{proof}
\end{homeworkProblem}

\newpage

\begin{homeworkProblem}
  Let $W(k)$ be the smallest integer $n$ such that any two-colouring of the set $\{1, 2, \dots, n\}$ contains a monochromatic arithmetic progression of length $k$.
  
  \begin{enumerate}[(a)]
    \item Use the first-moment method to show that $W(k) \ge 2^{k/2}$.
    \begin{proof}
      Randomly two-color the set $\{1, 2, \dots, n\}$, where $n < 2^{k/2}$. Let $X$ be the number of monochromatic arithmetic progressions of length $k$. Note that the number of $k$-APs in $[n]$ is 
      \[
        \sum_{a = 1}^{n - 1}\frac{n - a}{k - 1} = \frac{n(n - 1)}{2(k - 1)} \leq \frac{n^2}{2(k - 1)}.
      \]
      But then
      \[
        \E[X] = \frac{n^2}{2(k - 1)} \cdot \frac{2}{2^k} < 1.
      \]
      Thus there exists a two-coloring of $\{1, 2, \dots, n\}$ that contains an arithmetic progression of length $k$ that is not monochromatic.
    \end{proof}
    \item Use the Local Lemma to show that $W(k) \ge \frac{2^k}{2ek}(1 + o(1))$.
    \begin{proof}
      Suppose $n < \frac{2^k}{2ek}(1 + o(1))$. Randomly two-color the set $\{1, 2, \dots, n\}$. For each arithmetic progression $A$, let $E_A$ be the event that $A$ is monochromatic. Note that $E_A$ is only dependent on the events $E_{A'}$ where $A'$ overlaps with $A$ by at least 1 element. Since each AP is determined by it's first element and the common difference, for each $a \in A$ there are at most $k \cdot \frac{n}{k - 1}$ APs that contain $a$. Thus there are at most $d = k^2 \cdot \frac{n}{k - 1} - 1$ APs that overlap with $A$ by at least 1 element. Since
      \[
        \prob(E_A) = 2^{1 - k} < \frac{k - 1}{e \cdot k \cdot \frac{2^k}{2ek}} \leq \frac{1}{ed},
      \]
      the result now follows from the Symmetric Local Lemma.
    \end{proof}
  \end{enumerate}
\end{homeworkProblem}

\end{document}